\begin{textsample}{POS Dim 2 – 2009\_c5 – Score -7.00 – 2009\_c5/enade\_2009\_ri\_4686cb254d.txt}  \label{ex:f2_pos_418}
Capítulo 5 Distribuição dos Conceitos 5.1 Panorama nacional da distribuição dos conceitos A Tabela 5.1 apresenta a quantidade e o percentual de cursos de RELAÇÕES INTERNACIONAIS participantes do ENADE/2009 que se situaram em cada uma das faixas de conceitos, distribuídos de acordo com a Região em que se localizam os cursos das instituições de educação superior, mantenedoras dos cursos responsáveis pela oferta.
O Gráfico 5.1 apresenta a mesma informação sob a forma de diagrama de caixa para uma melhor visualização da dispersão dos conceitos obtidos.
No diagrama de caixa, a caixa propriamente dita fica limitada pelos quartis.
A barra dentro da caixa em negrito representa a mediana.
Linhas verticais unem os quartis aos valores extremos.
Se estes valores extremos se configurarem como valores atípicos, estes aparecem como pontos isolados e as linhas terminam no maior ( ou menor ) valor ainda não considerado atípico.
Eventuais valores atípicos aparecem como pequenos círculos com a Categoria Administrativa correspondente.
Valores considerados muito atípicos foram representados por estrelas.
Um diagrama de caixa para poucos dados ou com dados muito concentrados pode ter alguns pontos coincidentes.
Por exemplo, se mais de metade dos cursos de uma Região receberem conceito 5, o valor extremo superior, o quartil superior e a mediana coincidirão.
Se o conjunto só contém um dado, este apareceria denotado como a barra da mediana.
Se o conjunto apresenta todas as notas iguais, a notação é a mesma, somente uma barra horizontal.
Observando -se os dados da Tabela 5.1, percebe -se que, em termos nacionais, o conceito modal foi o conceito 2, correspondendo a 25,9 \% do total nacional de cursos.
Foram classificados com conceito 1, 16,6 \% dos cursos.
Os conceitos 3 e 4 foram alocados a, respectivamente, 18,5 \% e 6,2 \%.
O conceito máximo foi obtido por 4,9 \% dos cursos.
Dos cursos participantes, a maioria, 30,9 \%, ficou sem conceito.
Considerando -se separadamente as cinco Regiões, observou -se que, da Região Norte, 2 cursos participaram, dos quais 1, ficou sem conceito e 1 curso obteve conceito 2.
Esta Região foi a que apresentou cursos com a pior distribuição de conceitos.
Na Região Nordeste, 7 cursos participaram, dos quais, 2 ficaram sem conceito, 1 obteve conceito 2 e 3 cursos, o conceito 3.
Nenhum curso desta Região obteve conceito máximo e 1 curso obteve o conceito mínimo.
Na Região Sudeste, 43 cursos participaram, dos quais, 11 ficaram sem conceito ( 25,6 \% ), 13 ( 30,2 \% ) obtiveram conceito 2, 6 ( 14,0 \% ), o conceito 3 e 4 cursos ( 9,3 \% ), o conceito 4.
Nesta Região 2 cursos ( 4,7 \% ) obtiveram conceito máximo e 7 cursos, o conceito mínimo ( 16,3 \% ).
Na Região Sul, 18 cursos participaram, dos quais, 9 ficaram sem conceito, 2 obtiveram conceito 2, 3 cursos, o conceito 3 e 1 curso, o conceito 4.
Nesta Região 1 curso ( 5,6 \% ) obteve conceito máximo e 2 cursos, o conceito mínimo ( 11,1 \% ).
Na Região Centro-Oeste, 11 cursos participaram, dos quais, 2 ficaram sem conceito ( 18,2 \% ), 4 obtiveram conceito 2 ( 36,4 \% ) e 3 cursos o conceito 3 ( 27,3 \% ).
Nesta Região 1 curso ( 9,1 \% ) obteve conceito máximo e 1 curso, o conceito mínimo ( 9,1 \% ). 5.2 Conceitos por Categoria Administrativa e por Região A Tabela 5.2 apresenta a distribuição dos cursos participantes do ENADE/2009 de RELAÇÕES INTERNACIONAIS, por Categoria Administrativa, de acordo com os conceitos por eles alcançados, segundo a Região.
O Gráfico 5.2 apresenta a mesma informação sob a forma de diagrama de caixa para uma melhor visualização da dispersão dos conceitos obtidos.
Somente eventuais valores atípicos aparecem com o identificador da Região na qual o curso foi habilitado.
De acordo com as informações apresentadas, a grande maioria, 85,2 \%, é de cursos em instituições privadas, 11,1 \% dos cursos eram vinculados a instituições federais e 3,7 \% a instituições estaduais.
Nesta área, não houve participação de cursos em instituições municipais.
As instituições federais foram as que apresentaram, em termos nacionais, maior proporção de cursos com conceito 5, equivalentes a 11,1 \% do total dessa categoria.
Dos 9 cursos em instituições federais, 7 ficaram sem conceito ( 77,8 \% ).
Apenas 1 curso recebeu conceito 4 e 1 curso, o conceito máximo.
Entre as instituições estaduais, dos 3 cursos participantes, 2 receberam conceito 3 e 1 conceito 4.
Dos 69 cursos nas instituições privadas, 18 ficaram sem conceito com uma incidência de 26,1 \%.
Somente 3 cursos ( 4,3 \% ) receberam o conceito máximo, 11 cursos ( 18,9 \% ) receberam o conceito mínimo, 21 cursos ( 30,4 \% ) ficaram com conceito 2 e 13 receberam conceito 3 ( 18,8 \% ).
Na Região Norte, dos 2 cursos participantes, 1 era de instituição federal e 1 de instituição privada.
O curso de instituição federal ficou sem conceito e o curso de instituição privada recebeu conceito 2.
Na Região Nordeste, dos 7 cursos participantes, 1 era de instituição federal, 1 de instituição estadual e 5 de instituições privadas.
O único curso de instituição federal ficou sem conceito.
O único curso de instituição estadual recebeu conceito 3.
Dos 5 cursos em instituições privadas, 1 ficou sem conceito, 1 recebeu conceito 2, 1 recebeu o conceito mínimo e 2 o conceito 3.
Na Região Sudeste, dos 43 cursos participantes, 3 eram de instituições federais, 2 de instituições estaduais e 38 de instituições privadas.
Os 3 cursos de instituições federais ficaram sem conceito.
Dos 2 cursos de instituições estaduais, 1 curso ficou com conceito 3 e 1 com conceito 4.
Dos 38 cursos em instituições privadas, 8 ficaram sem conceito, 13 receberam conceito 2, 5 o conceito 3 e 3 o conceito 4.
Nas instituições privadas desta Região, 7 cursos receberam conceito mínimo e 2 cursos, conceito máximo.
Na Região Sul dos 18 cursos participantes, 2 eram de instituições federais e 16 de instituições privadas.
Não houve participação de cursos em instituição estadual ou municipal, nesta Região.
Dos 2 cursos de instituições federais, 1 curso ficou sem conceito e o outro obteve o conceito 4.
Dos 16 cursos em instituições privadas, 8 ficaram sem conceito, 2 receberam conceito 2 e 3 o conceito 3.
Nas instituições privadas desta Região, 2 cursos receberam conceito mínimo e 1 curso, conceito máximo.
Na Região Centro-Oeste, dos 11 cursos participantes, 2 eram de instituições federais e 9 de instituições privadas.
Não houve participação de cursos de instituição estadual ou municipal, nesta Região.
Dos 2 cursos de instituições federais, 1 curso ficou sem conceito e o outro obteve o conceito máximo.
Dos 9 cursos em instituições privadas, 1 ficou sem conceito, 4 receberam conceito 2, 1 o conceito mínimo e 3 o conceito 3. 5.3 Conceitos por Organização Acadêmica e por Região Na Tabela 5.3 encontra -se a distribuição dos conceitos atribuídos aos cursos participantes do ENADE/2009 de RELAÇÕES INTERNACIONAIS, por Organização Acadêmica, segundo as Regiões.
O Gráfico 5.3 apresenta a mesma informação sob a forma de diagrama de caixa para uma melhor visualização da dispersão dos conceitos obtidos.
Somente eventuais valores atípicos apareceram com o identificador da Categoria Administrativa.
Os cursos desta área estão principalmente alocados em Universidades ( 42,0 \% ), depois em Faculdades ( 32,1 \% ), e Centros Universitários ( 25,9 \% ).
Nas Universidades e Centros Universitários, o conceito modal foi 2 ( 23,5 \% e 42,9 \%, respectivamente ).
Nas Faculdades, o conceito modal foi o conceito 3 ( 19,2 \% ).
Nos 3 tipos de Organização Acadêmica, Universidades, Centros Universitários e Faculdades, receberam conceito mínimo, respectivamente 11,8 \%, 23,8 \% e 7,7 \%.
Uma proporção menor de cursos atingiu o conceito máximo : 8,8 \% em Universidades e 3,8 \% em Faculdades.
Nenhum curso de Centro Universitário obteve conceito máximo.
Nos 3 tipos de Organização Acadêmica, Universidades, Centros Universitários e Faculdades, existiram cursos que ficaram sem conceito, respectivamente 32,4 \%, 4,8 \% e 50,0 \%.
Considerando -se separadamente as cinco Regiões, observou -se que, na Região Norte, dos 2 cursos participantes, 2 eram de Universidades, 1 ficou sem conceito e 1 obteve o conceito 2.
Na Região Nordeste, dos 7 cursos participantes, 3 eram de Universidades, 2 de Centros Universitários e 2 cursos de Faculdades.
Dos 3 cursos em Universidades, 2 ficaram sem conceito e 1 obteve o conceito 3.
Os 2 cursos em Centros Universitários obtiveram conceitos 1 e 3.
Dos 2 cursos em Faculdades, 1 obteve o conceito 2 e o outro, conceito 3.
Na Região Sudeste, dos 43 cursos participantes, 17 eram de Universidades, 12 de Centros Universitários e 14 cursos de Faculdades.
Dos 17 cursos em Universidades, 4 ficaram sem conceito, 5 obtiveram conceito 2, 1 curso, o conceito 3 e 2 cursos, o conceito 4.
Nas Universidades, 3 cursos obtiveram conceito mínimo e 2 cursos, o conceito máximo.
Dos 12 cursos em Centros Universitários, 6 obtiveram conceito 2, 1 o conceito 3 e 1 o conceito 4.
Nenhum curso em Centro Universitário obteve conceito máximo e 4 obtiveram conceito mínimo.
Dos 14 cursos em Faculdades, 7 ficaram sem conceito, 2 obtiveram o conceito 2, 4 cursos, conceito 3 e 1 curso o conceito 4.
Nenhum curso em Faculdades obteve conceito mínimo ou máximo.
Na Região Sul, dos 18 cursos participantes, 8 eram de Universidades, 3 de Centros Universitários e 7 cursos de Faculdades.
Dos 8 cursos em Universidades, 3 ficaram sem conceito, 2 obtiveram conceito 3, 1 curso, o conceito 2 e 1 curso, o conceito 4.
Nas Universidades, 1 curso obteve conceito mínimo e nenhum curso, o conceito máximo.
Dos 3 cursos em Centros Universitários, 1 ficou sem conceito, 1 obteve conceito 2 e 1 o conceito 3.
Nenhum curso em Centro Universitário obteve conceito máximo nem conceito mínimo.
Dos 7 cursos em Faculdades, 5 ficaram sem conceito, 1 obteve o conceito mínimo e 1 o conceito máximo.
Na Região Centro-Oeste, dos 11 cursos participantes, 4 eram de Universidades, 4 de Centros Universitários e 3 cursos de Faculdades.
Dos 4 cursos em Universidades, 1 ficou sem conceito, 1 obteve conceito 2, 1 curso, o conceito 3 e 1 cursos obteve conceito máximo.
Dos 4 cursos em Centros Universitários, 2 obtiveram conceito 2 e 2 cursos o conceito 3.
Nenhum curso em Centro Universitário obteve conceito máximo nem conceito mínimo.
Dos 3 cursos em Faculdades, 1 ficou sem conceito, 1 obteve o conceito 2 e 1 o conceito mínimo.

% matched lemmas: 
\end{textsample}
